\documentclass[12pt]{article}
\usepackage{amsmath, amssymb}
\title{S2 G3 ITS L9 Scribe CoT AU2440248 Aryan}
\author{Aryan Thakkar}
\date{\today}


\begin{document}
\maketitle
\section*{CSE400: Fundamentals of Probability in Computing}
\subsection*{Lecture 9: Uniform, Exponential, Laplace and Gamma Random Variables}
\subsubsection*{(As per lecture slides dated February 2, 2026)}

\section*{1. Outline of the Lecture}

Types of Continuous Random Variables

Uniform Random Variable

Example

Exponential Random Variable

Example

Laplace Random Variable

Example

Gamma Random Variable

Graph and Special Cases

Example

Homework Problem

Problem Solving

In-Class Activity: Gaussian Density Estimation

\section*{2. Types of Continuous Random Variables}

\section*{3. Uniform Random Variable}

\subsection*{Definition and Notation}

Let $X$ be a uniform random variable on the interval $[a,b]$.

\subsection*{Probability Density Function (PDF)}

\[
f_X(x) =
\begin{cases}
\dfrac{1}{b-a}, & a \le x < b, \\
0, & \text{elsewhere}.
\end{cases}
\]

\subsection*{Cumulative Distribution Function (CDF)}

\[
F_X(x) =
\begin{cases}
0, & x < a, \\
\dfrac{x-a}{b-a}, & a \le x < b, \\
1, & x \ge b.
\end{cases}
\]

The PDF is constant over $[a,b]$, and the CDF increases linearly from $0$ to $1$ over the same interval.

\subsection*{Graphical Representation}

(a) Probability Density Function of a uniform random variable

(b) Cumulative Distribution Function of a uniform random variable

(Figure 3.8 in the lecture slides)

\section*{4. Example \#1: Uniform Random Variable}

\subsection*{Problem Statement}

The phase of a sinusoid, $\Theta$, is uniformly distributed over $[0,2\pi]$ such that its PDF is:

\[
f_{\Theta}(\theta) =
\begin{cases}
\dfrac{1}{2\pi}, & 0 \le \theta < 2\pi, \\
0, & \text{otherwise}.
\end{cases}
\]

\subsection*{(a) Find $\Pr(\Theta > \frac{3\pi}{4})$}

For a uniform random variable on $[0,2\pi]$:

\[
\Pr(a < \Theta < b) = \dfrac{b-a}{2\pi}
\]

\[
\Pr\left(\Theta > \frac{3\pi}{4}\right)
= \dfrac{2\pi - \frac{3\pi}{4}}{2\pi}
= \dfrac{5}{8}
\]

\subsection*{(b) Find $\Pr(\Theta < \pi \mid \Theta > \frac{3\pi}{4})$}

Using conditional probability:

\[
\Pr(A \mid B) = \dfrac{\Pr(A \cap B)}{\Pr(B)}
\]

\[
\Pr\left(\frac{3\pi}{4} < \Theta < \pi\right)
= \dfrac{\pi - \frac{3\pi}{4}}{2\pi}
= \dfrac{1}{8}
\]

\[
\Pr\left(\Theta > \frac{3\pi}{4}\right) = \dfrac{5}{8}
\]

\[
\Pr\left(\Theta < \pi \mid \Theta > \frac{3\pi}{4}\right)
= \dfrac{1/8}{5/8}
= \dfrac{1}{5}
\]

\subsection*{(c) Find $\Pr(\cos \Theta < \frac{1}{2})$}

\[
\cos \Theta = \frac{1}{2}
\Rightarrow \Theta = \frac{\pi}{3}, \frac{5\pi}{3}
\]

\[
\cos \Theta < \frac{1}{2}
\quad \text{for} \quad
\frac{\pi}{3} < \Theta < \frac{5\pi}{3}
\]

\[
\Pr(\cos \Theta < \tfrac{1}{2})
= \dfrac{\frac{5\pi}{3} - \frac{\pi}{3}}{2\pi}
= \dfrac{4\pi/3}{2\pi}
= \dfrac{2}{3}
\]

\section*{5. Applications of Uniform Random Variable}

Phase of a sinusoidal signal when all phase angles between $0$ and $2\pi$ are equally likely

A random number generated by a computer between $0$ and $1$ for simulations

Arrival time of a user within a known time window, assuming no time preference

\section*{6. Exponential Random Variable}

\subsection*{Definition}

The exponential random variable has the following PDF and CDF (for any $b > 0$):

\subsection*{Probability Density Function}

\[
f_X(x) = \frac{1}{b} e^{-x/b} u(x)
\]

\subsection*{Cumulative Distribution Function}

\[
F_X(x) = \left[1 - e^{-x/b}\right] u(x)
\]

where $u(x)$ is the unit step function.

\subsection*{Graphical Representation}

(a) PDF of exponential random variable

(b) CDF of exponential random variable for $b = 2$

\section*{7. Example \#2: Exponential Random Variable}

\subsection*{Problem Statement}

Let $X$ be an exponential random variable with PDF:

\[
f_X(x) = e^{-x} u(x)
\]

\subsection*{(a) Find $\Pr(3X < 5)$}

\[
\Pr(3X < 5) = \Pr\left(X < \frac{5}{3}\right)
\]

Using the CDF:

\[
F_X(x) = 1 - e^{-x}
\]

\[
\Pr\left(X < \frac{5}{3}\right)
= 1 - e^{-5/3}
\]

\subsection*{(b) Generalize your answer to find $\Pr(3X < y)$}

\[
\Pr(3X < y) = \Pr\left(X < \frac{y}{3}\right)
= 1 - e^{-y/3}
\]

\section*{8. Remaining Topics Mentioned in Outline}

Laplace Random Variable

Gamma Random Variable

Graph and special cases

Example

Homework problem

(These topics are listed in the lecture outline but are not developed further in the provided slides of this PDF.)

\section*{End of Lecture 9 Scribe}

\end{document}
